% Part A. Included from discussion.tex with % Part A. Included from discussion.tex with % Part A. Included from discussion.tex with % Part A. Included from discussion.tex with \input{sections/partA_discussion}.
\subsection{Preprocessing and noise level}
\label{sec:partA-discussion}
\subsubsection{The filters lower the peaks a little}
The median and Savitzky--Golay filters remove noise, but they also trim the top of narrow
peaks: the 20 tallest peaks of the Big pieces recording keep a median of 96\,\% of their height,
and the worst one only 59\,\%. The heights and areas used for the amount are therefore a little
too low for narrow peaks. A shorter Savitzky--Golay window, or no median filter, would keep the
peaks but lets more noise through, and we did not test which setting is best for the amount.

\subsubsection{One setting, one noise level per recording}
The noise level $\sigma$ is measured from each recording, which is needed: it differs by a factor
of 8 (0.51--4.03 counts), and one common value changes the piece count in 23 of 28 recordings.
However, $\sigma$ is only measured correctly if the recording has quiet parts. In a
recording with flow all the time, the quietest blocks still contain signal, $\sigma$ is too high
and small pieces are lost. The factor $k = 10$ is the same for all recordings. The piece count
changes with $k$ (Table~\ref{tab:k-sweep}), and with no known number of pieces in the dense
recordings we cannot say which $k$ is right. The recording with the highest $\sigma$ (4.03, a
tumour recording) shows that a high noise level is possible and should be checked by eye.

\subsubsection{What would fix them}
A recording with known pieces and several flow rates would show which $k$ gives the right count.
A noise level measured from a short recording without suction at the start of each session would
not depend on quiet blocks inside the suction.
.
\subsection{Preprocessing and noise level}
\label{sec:partA-discussion}
\subsubsection{The filters lower the peaks a little}
The median and Savitzky--Golay filters remove noise, but they also trim the top of narrow
peaks: the 20 tallest peaks of the Big pieces recording keep a median of 96\,\% of their height,
and the worst one only 59\,\%. The heights and areas used for the amount are therefore a little
too low for narrow peaks. A shorter Savitzky--Golay window, or no median filter, would keep the
peaks but lets more noise through, and we did not test which setting is best for the amount.

\subsubsection{One setting, one noise level per recording}
The noise level $\sigma$ is measured from each recording, which is needed: it differs by a factor
of 8 (0.51--4.03 counts), and one common value changes the piece count in 23 of 28 recordings.
However, $\sigma$ is only measured correctly if the recording has quiet parts. In a
recording with flow all the time, the quietest blocks still contain signal, $\sigma$ is too high
and small pieces are lost. The factor $k = 10$ is the same for all recordings. The piece count
changes with $k$ (Table~\ref{tab:k-sweep}), and with no known number of pieces in the dense
recordings we cannot say which $k$ is right. The recording with the highest $\sigma$ (4.03, a
tumour recording) shows that a high noise level is possible and should be checked by eye.

\subsubsection{What would fix them}
A recording with known pieces and several flow rates would show which $k$ gives the right count.
A noise level measured from a short recording without suction at the start of each session would
not depend on quiet blocks inside the suction.
.
\subsection{Preprocessing and noise level}
\label{sec:partA-discussion}
\subsubsection{The filters lower the peaks a little}
The median and Savitzky--Golay filters remove noise, but they also trim the top of narrow
peaks: the 20 tallest peaks of the Big pieces recording keep a median of 96\,\% of their height,
and the worst one only 59\,\%. The heights and areas used for the amount are therefore a little
too low for narrow peaks. A shorter Savitzky--Golay window, or no median filter, would keep the
peaks but lets more noise through, and we did not test which setting is best for the amount.

\subsubsection{One setting, one noise level per recording}
The noise level $\sigma$ is measured from each recording, which is needed: it differs by a factor
of 8 (0.51--4.03 counts), and one common value changes the piece count in 23 of 28 recordings.
However, $\sigma$ is only measured correctly if the recording has quiet parts. In a
recording with flow all the time, the quietest blocks still contain signal, $\sigma$ is too high
and small pieces are lost. The factor $k = 10$ is the same for all recordings. The piece count
changes with $k$ (Table~\ref{tab:k-sweep}), and with no known number of pieces in the dense
recordings we cannot say which $k$ is right. The recording with the highest $\sigma$ (4.03, a
tumour recording) shows that a high noise level is possible and should be checked by eye.

\subsubsection{What would fix them}
A recording with known pieces and several flow rates would show which $k$ gives the right count.
A noise level measured from a short recording without suction at the start of each session would
not depend on quiet blocks inside the suction.
.
\subsection{Preprocessing and noise level}
\label{sec:partA-discussion}
\subsubsection{The filters lower the peaks a little}
The median and Savitzky--Golay filters remove noise, but they also trim the top of narrow
peaks: the 20 tallest peaks of the Big pieces recording keep a median of 96\,\% of their height,
and the worst one only 59\,\%. The heights and areas used for the amount are therefore a little
too low for narrow peaks. A shorter Savitzky--Golay window, or no median filter, would keep the
peaks but lets more noise through, and we did not test which setting is best for the amount.

\subsubsection{One setting, one noise level per recording}
The noise level $\sigma$ is measured from each recording, which is needed: it differs by a factor
of 8 (0.51--4.03 counts), and one common value changes the piece count in 23 of 28 recordings.
However, $\sigma$ is only measured correctly if the recording has quiet parts. In a
recording with flow all the time, the quietest blocks still contain signal, $\sigma$ is too high
and small pieces are lost. The factor $k = 10$ is the same for all recordings. The piece count
changes with $k$ (Table~\ref{tab:k-sweep}), and with no known number of pieces in the dense
recordings we cannot say which $k$ is right. The recording with the highest $\sigma$ (4.03, a
tumour recording) shows that a high noise level is possible and should be checked by eye.

\subsubsection{What would fix them}
A recording with known pieces and several flow rates would show which $k$ gives the right count.
A noise level measured from a short recording without suction at the start of each session would
not depend on quiet blocks inside the suction.
